% TEMPLATE-MILC-v3.tex - plantilla maestra UNICA de las guias MILC v3.
%
% La usan las tres familias de guias, cada una con su driver:
%   · Grados 6-11  -> scripts/build-guias-g11.py
%   · Bebras       -> scripts/build-guias-bebras.py
%   · Semillero    -> scripts/build-guias-semillero.py
%
% Antes habia tres copias casi identicas (~400 lineas iguales, 6 distintas):
% cada mejora habia que aplicarla tres veces, y en la practica no se hacia
% (el motor de recursos vivia solo en dos de ellas). Lo que varia entre
% familias esta parametrizado en 6 placeholders que cada driver llena
% (se nombran aqui SIN su sintaxis real: el builder reemplaza por texto plano,
% tambien dentro de los comentarios, y un valor multilinea rompe el preambulo):
%
%   HEADER_IZQ        encabezado izquierdo de pagina
%   HEADER_DER        encabezado derecho de pagina
%   PORTADA_ETIQUETA  rotulo sobre el numero grande (GRADO/MOMENTO/MODULO)
%   PORTADA_SERIE     linea de serie de la portada
%   PORTADA_ID        identificador de la guia en la portada
%   PORTADA_CIRCULO   el circulito de grado (°), o vacio si no aplica
%
% El resto de claves las documenta content/guias/_SCHEMA.md. El script valida
% que no queden placeholders sin reemplazar antes de escribir el archivo final.

\documentclass[11pt,letterpaper]{article}
\usepackage[margin=2.0cm]{geometry}
\usepackage[table]{xcolor}
\usepackage{fontspec}
\usepackage[spanish]{babel}
\usepackage{tikz}
% Librerías TikZ para los diagramas de las guías (bloque `recursos:`):
%  · babel  → babel en español vuelve activos caracteres como «>», y TikZ los usa
%             en sus opciones (>=stealth, ->). Sin esta librería, cualquier
%             diagrama con flechas revienta con «\language@active@arg> has an
%             extra }». Debe ir ANTES de que se dibuje nada.
%  · positioning        → sintaxis «below left=2pt and 3pt».
%  · decorations.pathreplacing → llaves ({brace}) para señalar rangos.
\usetikzlibrary{babel,positioning,decorations.pathreplacing}
\usepackage{graphicx}
% Circuitos: el trabajo de electrónica se hace hoy con micro:bit y sus sensores
% integrados (MakeCode, sin cableado), así que no se necesitan esquemáticos.
% Si algún día entran montajes con protoboard, basta descomentar esta línea:
% los diagramas .tex/.tikz declarados en `recursos:` ya se incrustan solos.
% \usepackage{circuitikz}
\usepackage[most]{tcolorbox}
\usepackage{tabularx}
\usepackage{array}
\usepackage{fancyhdr}
\usepackage{titlesec}
\usepackage[document]{ragged2e}  % bandera a la derecha en todo el cuerpo
\usepackage{enumitem}
\usepackage{microtype}

% Helvetica Neue no trae la flecha U+2192, asi que cada «→» escrito literalmente
% en el YAML se compilaba a NADA, en silencio (462 flechas en 61 guias: rutas del
% tipo «paleta → tipografia → lockup» salian rotas). Mapeamos el caracter a su
% version matematica, que si tiene glifo. Asi el autor puede seguir escribiendo
% «→» comodamente en el YAML.
\usepackage{newunicodechar}
\newunicodechar{→}{\ensuremath{\rightarrow}}
% Mismo problema con los simbolos de marca y vinetas: el «✓» de «una palomita ✓»
% salia como cuadrito vacio en el PDF de todas las guias que lo usaban.
\usepackage{pifont}
\newunicodechar{✓}{\ding{51}}
\newunicodechar{✗}{\ding{55}}
\newunicodechar{✔}{\ding{52}}
\newunicodechar{•}{\textbullet}
\newunicodechar{≥}{\ensuremath{\geq}}
\newunicodechar{≤}{\ensuremath{\leq}}

\setmainfont{Helvetica Neue}
\setsansfont{Helvetica Neue}
\setlength{\parindent}{0pt}
\setlength{\parskip}{6pt}
% Sin particion de palabras: en 6.º-7.º una palabra cortada por guion al final
% de linea («Comuni-cacion») frena la lectura mucho mas que una linea corta.
\hyphenpenalty=10000
\exhyphenpenalty=10000
\pretolerance=2000
\tolerance=3000
\setlength{\emergencystretch}{3em}
\linespread{1.10}
\renewcommand{\arraystretch}{1.25}
\setlist{leftmargin=5mm,itemsep=1mm,topsep=1mm}
\newcolumntype{Y}{>{\RaggedRight\arraybackslash}X}

\definecolor{milcMagenta}{HTML}{E6007E}
\definecolor{milcVino}{HTML}{5A0038}
\definecolor{milcTurquesa}{HTML}{0093A5}
\definecolor{milcVerde}{HTML}{58A923}
\definecolor{milcMostaza}{HTML}{E5B400}
\definecolor{milcOcre}{HTML}{B86600}
\definecolor{milcNegro}{HTML}{241816}
\definecolor{milcGris}{HTML}{F7F7F4}
\definecolor{milcLinea}{HTML}{E8E2DD}
\definecolor{milcGradePrimary}{HTML}{0D9488}
\definecolor{milcGradeDark}{HTML}{115E59}
\definecolor{milcGradeSoft}{HTML}{CCFBF1}
\definecolor{milcGradeAccent}{HTML}{CFE7FF}
\definecolor{milcGradeText}{HTML}{FFFFFF}
\color{milcNegro}

\pagestyle{fancy}
\fancyhf{}
\lhead{\small Territorio interior · Educación socioemocional}
\rhead{\small Grado 8° · Momento 2 · MILC v3}
\cfoot{\small\thepage}
\renewcommand{\headrulewidth}{0pt}

\newtcolorbox{titlebox}[1]{
  enhanced,colback=#1,colframe=#1,arc=7mm,boxrule=0pt,
  left=7mm,right=7mm,top=3mm,bottom=3mm,
  before skip=8pt,after skip=8pt,
  fontupper=\sffamily\bfseries\Large\color{white}
}

\newtcolorbox{softbox}[3]{
  enhanced,colback=#2,colframe=#1,borderline west={4pt}{0pt}{#1},
  arc=2mm,boxrule=.4pt,left=4mm,right=4mm,top=3mm,bottom=3mm,
  before skip=6pt,after skip=8pt,title={#3},coltitle=white,
  colbacktitle=#1,fonttitle=\sffamily\bfseries,fontupper=\RaggedRight,
  attach boxed title to top left={xshift=3mm,yshift=-2mm},
  boxed title style={arc=2mm, boxrule=0pt}
}

\newtcolorbox{drawbox}[2]{
  enhanced,colback=white,colframe=#1,arc=4mm,boxrule=1pt,
  left=5mm,right=5mm,top=4mm,bottom=4mm,before skip=8pt,after skip=8pt,
  title={#2},coltitle=white,colbacktitle=#1,fonttitle=\sffamily\bfseries,
  fontupper=\RaggedRight,
  attach boxed title to top left={xshift=4mm,yshift=-2mm},
  boxed title style={arc=3mm, boxrule=0pt}
}

\newtcolorbox{infoband}[1]{
  enhanced,colback=#1!8,colframe=#1!50,arc=2mm,boxrule=.5pt,
  left=4mm,right=4mm,top=2mm,bottom=2mm,before skip=4pt,after skip=4pt,
  fontupper=\small\RaggedRight
}

% Puente narrativo entre secciones (contrato editorial Conéctate).
% Da continuidad: "Acabas de X. Ahora vas a Y porque Z."
% Visualmente: banda izquierda en color del grado, texto en itálica pequeña.
\newtcolorbox{puentebox}{
  enhanced,colback=milcGradeSoft,colframe=milcGradeDark,
  borderline west={3pt}{0pt}{milcGradeDark},
  arc=0mm,boxrule=0pt,left=5mm,right=4mm,top=2.5mm,bottom=2.5mm,
  before skip=4pt,after skip=8pt,
  fontupper=\itshape\small\color{milcNegro}\RaggedRight
}
% La flecha va en modo matematico a proposito: Helvetica Neue NO tiene el glifo
% U+2192, asi que \textrightarrow se compilaba a NADA («Missing character: There
% is no → in font Helvetica Neue») y el puente salia sin flecha. En modo
% matematico la flecha viene de la fuente matematica y si se dibuja.
\newcommand{\puente}[1]{%
  \begin{puentebox}%
    \textcolor{milcGradeDark}{$\rightarrow$}\hspace{2mm}#1%
  \end{puentebox}%
}

\newcommand{\linead}[1]{\noindent\color{milcLinea}\rule{#1}{0.5pt}\color{milcNegro}}
\newcommand{\respuesta}[1]{\par\vspace{2mm}\foreach \n in {1,...,#1}{\noindent\color{milcLinea}\rule{\linewidth}{0.5pt}\par\vspace{5mm}}\color{milcNegro}}
\newcommand{\checkbox}{\textcolor{milcGradeDark}{\fbox{\phantom{x}}}\hspace{5pt}}

% ─── Listas aireadas para las guías socioemocionales ────────────────────
% Componer las verificaciones como «\checkbox texto\\» deja el texto que salta
% de línea debajo del cuadrito y apelmaza el bloque. Estos entornos alinean la
% continuación y airean los ítems. Los usa Territorio Interior; cualquier
% familia puede hacerlo.
\newlist{checklista}{itemize}{1}
\setlist[checklista]{
  label=\textcolor{milcGradeDark}{\fbox{\phantom{x}}},
  leftmargin=9mm, labelsep=3.2mm, itemsep=2.4mm,
  topsep=2.5mm, parsep=0pt, partopsep=0pt, before=\RaggedRight
}
\newlist{pasoslista}{enumerate}{1}
\setlist[pasoslista]{
  label=\textcolor{milcGradeDark}{\sffamily\bfseries\arabic*.},
  leftmargin=8mm, labelsep=3mm, itemsep=2.8mm,
  topsep=1mm, parsep=0pt, partopsep=0pt, before=\RaggedRight
}

\newcommand{\phasecard}[4]{
\begin{tcolorbox}[enhanced,colback=#3,colframe=#2,arc=3mm,boxrule=.5pt,height=3.05cm,valign=center,halign=center,left=1.5mm,right=1.5mm,top=2mm,bottom=2mm]
{\sffamily\bfseries\fontsize{22}{24}\selectfont\textcolor{#2}{#1}}\par
{\sffamily\bfseries\small #4}
\end{tcolorbox}
}

% ── Piezas del rediseno 2026 (legibilidad 6.º-11.º) ───────────────────
% Banda de estacion: reemplaza el titlebox macizo. Titulo a la izquierda,
% dato practico (tiempo/modalidad) a la derecha, en una sola linea.
\newcommand{\milcestacion}[2]{%
\begin{tcolorbox}[enhanced,colback=#1,colframe=#1,arc=2mm,boxrule=0pt,
  left=5mm,right=5mm,top=2.6mm,bottom=2.6mm,before skip=11pt,after skip=7pt]
{\sffamily\bfseries\fontsize{15}{18}\selectfont\color{white}#2}
\end{tcolorbox}}

% Tarjeta de paso numerada: sustituye las tablas «Pilar / Aplicacion», que
% eran formato de consulta docente, no de estudiante. El numero va en un
% circulo sobre el borde para que la secuencia se lea de un vistazo.
\newcommand{\milcpaso}[3]{%
\begin{tcolorbox}[enhanced,colback=white,colframe=milcLinea,arc=1.5mm,boxrule=.9pt,
  left=14mm,right=4mm,top=3mm,bottom=3mm,before skip=4pt,after skip=4pt,
  fontupper=\RaggedRight,
  overlay={\node[anchor=north west,circle,fill=#2,text=white,inner sep=0pt,
    minimum size=7.4mm,font=\sffamily\bfseries\small]
    at ([xshift=3.6mm,yshift=-3.2mm]frame.north west) {#1};}]
#3
\end{tcolorbox}}

% Casilla marcable de tamano util con lapiz (la anterior era \fbox{\phantom{x}},
% de alto variable segun el texto que la seguia).
\newcommand{\milcmarca}{\framebox[3.8mm]{\rule{0pt}{3.4mm}}}

% Fila marcable de lista suelta (checks de la estacion 1).
\newcommand{\milccheckrow}[1]{%
\par\vspace{1.8mm}\noindent
\makebox[6.4mm][l]{\raisebox{-.55mm}{\textcolor{milcNegro}{\milcmarca}}}%
\parbox[t]{\dimexpr\linewidth-6.4mm\relax}{\RaggedRight #1}\par}

% Celda marcable dentro de la rubrica (tabla de autoevaluacion). Misma
% sangria francesa, pero pensada para vivir dentro de una columna X.
\newcommand{\milcopcion}[1]{%
\makebox[5.2mm][l]{\raisebox{-.55mm}{\textcolor{milcNegro}{\milcmarca}}}%
\parbox[t]{\dimexpr\linewidth-5.2mm\relax}{\RaggedRight #1}}

% ── Recursos visuales de la guía ──────────────────────────────────────
% Declarados en el bloque `recursos:` del YAML y emitidos por el builder.
% \guiaFigura   → imagen rasterizada o PDF (foto, carta celeste, montaje).
% \guiaDiagrama → fuente TikZ/circuitikz (.tex) incrustada como vector:
%                 es la vía para circuitos y esquemáticos editables.
% #1 = ruta relativa al .tex · #2 = pie (puede ir vacío)
\newcommand{\guiaFigura}[2]{%
\begin{center}
\includegraphics[width=0.88\linewidth,height=0.38\textheight,keepaspectratio]{#1}\\[1.5mm]
{\footnotesize\itshape\textcolor{milcNegro!75}{#2}}
\end{center}
}

\newcommand{\guiaDiagrama}[2]{%
\begin{center}
\input{#1}\\[1.5mm]
{\footnotesize\itshape\textcolor{milcNegro!75}{#2}}
\end{center}
}

\begin{document}
\thispagestyle{empty}

% ── Portada ───────────────────────────────────────────────────────────
% Antes: un rectangulo de color a pagina completa. Se veia bien en pantalla,
% pero en la fotocopiadora del colegio se comia el toner y salia gris sucio.
% Ahora la banda ocupa el tercio superior y el resto va en blanco.
\begin{tikzpicture}[remember picture,overlay]
  \path[fill=milcGradePrimary]
    (current page.north west) rectangle ([yshift=-10.6cm]current page.north east);

  \node[anchor=north west,text=milcGradeText] at ([xshift=2.0cm,yshift=-1.55cm]current page.north west)
    {\sffamily\bfseries\fontsize{12}{14}\selectfont MOMENTO};
  \node[anchor=north west,text=milcGradeText,opacity=.85] at ([xshift=2.0cm,yshift=-2.35cm]current page.north west)
    {\sffamily\bfseries\fontsize{8.5}{10}\selectfont TERRITORIO INTERIOR · SOCIOEMOCIONAL};
  \node[anchor=north east,text=milcGradeText,opacity=.85,align=right] at ([xshift=-2.0cm,yshift=-1.55cm]current page.north east)
    {\sffamily\bfseries\fontsize{9}{12}\selectfont I.E. Sor María Juliana\\Cartago, Valle del Cauca};

  \node[anchor=west,text=milcGradeText] (gradonum) at ([xshift=1.85cm,yshift=-7.1cm]current page.north west)
    {\sffamily\bfseries\fontsize{112}{112}\selectfont 2};
  % El circulito de grado (°) solo aplica a las guias de grado («11°»); en
  % Bebras y Semillero el numero grande es un momento/modulo, no un grado.
  % Se posiciona relativo al nodo `gradonum`, no en coordenadas absolutas.
  

  \node[anchor=west,text=milcGradeText,text width=11.4cm,align=left] at ([xshift=1.5cm,yshift=0.5cm]gradonum.east)
    {
     {\sffamily\bfseries\fontsize{9}{11}\selectfont\MakeUppercase{Territorio interior · Socioemocional}}\\[3mm]
     {\sffamily\bfseries\fontsize{21}{25}\selectfont\setlength{\baselineskip}{27pt}Decir\\[4pt]que no\\[4pt]negarse sin quedar por fuera}\\[3.5mm]
     {\sffamily\bfseries\fontsize{15}{18}\selectfont Grado 8° · Momento 2}
    };
\end{tikzpicture}

\vspace*{9.4cm}
\vfill

{\sffamily\bfseries\fontsize{8}{10}\selectfont\textcolor{milcGradeDark}{LO QUE TE LLEVAS HOY}}

\begin{tcolorbox}[enhanced,colback=white,colframe=milcNegro,arc=2mm,boxrule=1.6pt,
  left=5mm,right=5mm,top=4mm,bottom=4mm,before skip=3pt,after skip=12pt,
  fontupper=\RaggedRight\fontsize{11}{15.5}\selectfont]
Tu repertorio del no: tres negativas escritas y ensayadas —la corta, la que ofrece otra cosa y la que se apoya en un tercero—, más el nombre de la persona a la que le vas a pedir que sea tu aliado, y la petición ya hecha.
\end{tcolorbox}

\begin{tcolorbox}[enhanced,colback=milcGris,colframe=milcGris,arc=2mm,boxrule=0pt,
  left=5mm,right=5mm,top=3.5mm,bottom=3.5mm,after skip=12pt]
{\sffamily\bfseries\fontsize{8}{10}\selectfont\textcolor{milcNegro!55}{ESTA GUÍA ES DE}}\\[3mm]
\begin{tabularx}{\linewidth}{X p{3.2cm} p{3.2cm}}
\linead{\linewidth} & \linead{3.0cm} & \linead{3.0cm}\\[-1mm]
{\fontsize{8.5}{10}\selectfont\textcolor{milcNegro!55}{Nombre completo}} &
{\fontsize{8.5}{10}\selectfont\textcolor{milcNegro!55}{Grupo}} &
{\fontsize{8.5}{10}\selectfont\textcolor{milcNegro!55}{Fecha}}\\
\end{tabularx}
\end{tcolorbox}

\vfill

\noindent\textcolor{milcLinea}{\rule{\linewidth}{1pt}}\\[2mm]
{\sffamily\bfseries\fontsize{9.5}{12}\selectfont Autor: Dr. Álvaro Cárdenas Orozco}\\[1mm]
{\sffamily\fontsize{8.5}{11}\selectfont\textcolor{milcNegro!55}{Metodología MILC · Apertura ancestral · Presión de grupo · Triángulo Dussel-Estoicismo-Floridi}}

\newpage

\section*{Apertura: Decir que no: negarse sin quedar por fuera}

\begin{softbox}{milcGradeDark}{milcGradeSoft}{Saber ancestral}
El pueblo \textbf{misak}, en Guambía (Cauca), tiene un lugar donde se toma la palabra: el \textbf{Nu Nakchak}, el fogón. \emph{Ahí hablan los \textbf{shures} y las \textbf{shuras} ---los mayores---,} y de ahí sale lo que la comunidad decide. \textbf{Y de ese fogón salió uno de los «no» más grandes que ha dicho un pueblo en Colombia:} la \textbf{Misak Ley}. Los misak sostienen que tienen un \textbf{Derecho Mayor} ---anterior, propio, heredado de sus mayores por tradición oral--- y con él fundamentan su autoridad y su gobierno propio; el Cabildo de Guambía está compuesto por \textbf{tatas y mamas}, autoridades zonales y generales. \textbf{Fíjate en las tres cosas que hacen que ese «no» funcione.} \textbf{Primera: no lo dijo uno solo.} Salió del fogón, hablado entre mayores, no de una persona plantada de frente. \textbf{Segunda: dice a qué se dice \emph{sí}.} No es solo un rechazo ---afirma un derecho propio y una forma de gobernarse---. \textbf{Tercera: quedó por escrito}, así que no hay que volver a discutirlo cada vez. \emph{Compáralo con el «no» que uno intenta a los trece años: solo, de improviso, sin respaldo y sin nada que ofrecer a cambio. No es que a ti te falte carácter: es que ese «no» está armado para fallar.}
\end{softbox}

\begin{softbox}{milcTurquesa}{milcGris}{Saber contemporáneo}
¿Qué tanto puede el grupo? \textbf{Se midió, y el resultado es incómodo.} En los años cincuenta, \textbf{Solomon Asch} puso a una persona en un salón con otras que parecían participantes y eran actores. La tarea era ridículamente fácil: decir \textbf{cuál de tres rayas era igual de larga} a una de muestra. Los actores, uno tras otro, daban en voz alta \textbf{la respuesta equivocada}. \textbf{Y alrededor de un tercio de las veces, la persona real repetía el error} que estaba viendo con sus propios ojos. \textbf{Pero el hallazgo que de verdad importa es el otro.} Cuando Asch metió \textbf{una sola persona más} que contestara distinto del grupo ---\textbf{un aliado}---, \textbf{la conformidad se desplomó}: de cerca de un tercio a apenas unos pocos casos. \emph{Y funcionaba incluso cuando ese aliado también se equivocaba, pero en otra dirección.} \textbf{Traducido:} lo que rompe la presión \textbf{no es tener más carácter}: es \textbf{que la fila no sea unánime}. \emph{Basta con que una persona más diga otra cosa para que las demás se atrevan a confiar en lo que están viendo.}
\end{softbox}

\begin{softbox}{milcMostaza}{milcGris}{Pregunta-puente}
\textit{El «no» misak salió \textbf{de un fogón, entre varios, y por escrito}. \textbf{¿Cuál de tus «sí» de este año habrías podido convertir en un «no»} si hubieras tenido una sola persona al lado diciendo lo mismo que tú?}
\end{softbox}

\begin{softbox}{milcVerde}{milcGris}{Saber y saber hacer}
Al terminar podrás: (1) \textbf{explicar} por qué la presión de grupo funciona incluso cuando uno ve claramente lo contrario, y por qué casi nadie cree lo que está diciendo; (2) \textbf{reconocer} que un solo aliado desarma la unanimidad, que es de donde sale la presión; (3) \textbf{crear} tres negativas propias ---la corta, la que ofrece otra cosa y la que se apoya en un tercero---; (4) \textbf{pedirle} a alguien concreto que sea tu aliado, y serlo tú del otro.
\end{softbox}



\small
\begin{tabularx}{\textwidth}{>{\bfseries}p{3.0cm}Y}
\rowcolor{milcGradePrimary!12}Autor & Dr. Álvaro Cárdenas Orozco\\
DBA & Comprende los fenómenos que afectan a su territorio, regula sus emociones ante ellos y acompaña a otros con empatía (transversal 6°--11°, articulado con la política «Escuela, territorio de vida» del MEN, Resolución 006519 de 2025, y la Ley 1620 de 2013)\\
Referentes & Primera ayuda psicológica (OMS, 2011/2012) · Directrices IASC (2007) · USGS y Servicio Geológico Colombiano · Pueblo nasa y la minga (CRIC) · Dussel · Estoicismo · Floridi\\
Producto final & Tu repertorio del no: tres negativas escritas y ensayadas —la corta, la que ofrece otra cosa y la que se apoya en un tercero—, más el nombre de la persona a la que le vas a pedir que sea tu aliado, y la petición ya hecha.\\
\end{tabularx}
\normalsize

\puente{En el momento 1 viste \textbf{qué se paga por entrar}. Hoy toca lo siguiente: \textbf{cómo se deja de pagar} sin que te cueste el puesto.}

\milcestacion{milcGradeDark}{Ruta MILC: así vamos a aprender}
\begin{tabularx}{\textwidth}{>{\centering\arraybackslash}X>{\centering\arraybackslash}X>{\centering\arraybackslash}X>{\centering\arraybackslash}X}
\phasecard{1}{milcVerde}{milcGris}{Escucha\\[-1mm]\small Observo, escucho y pregunto.} &
\phasecard{2}{milcTurquesa}{milcGris}{Sistematización\\[-1mm]\small Me niego y me quedo.} &
\phasecard{3}{milcMagenta}{milcGris}{Praxis\\[-1mm]\small Construyo y aplico.} &
\phasecard{4}{milcVino}{milcGris}{Evaluación\\[-1mm]\small Reflexión liberadora.}
\end{tabularx}


\puente{Primero, tus propios síes: los que dijiste este año y no querías decir. \emph{Sin juzgarte: el experimento de Asch dice que le pasa a todo el mundo.}}

\milcestacion{milcVerde}{Estación 1 de 3 · Escucha}

\begin{softbox}{milcVerde}{milcGris}{Escena para escuchar}
\textbf{Actividad 1 · IDENTIFICA --- Los síes que no eran.} \textbf{Nadie va a leer esta página.} \textbf{Primera parte.} Escribe \textbf{tres veces de este año en que dijiste que sí y no querías}: qué era, quién lo pidió, cuánta gente había mirando. \emph{Pueden ser cosas mínimas ---prestar algo, ir a un lado, reírte de algo, cubrir a alguien---: las mínimas son las que más enseñan.} \textbf{Segunda parte.} Al frente de cada una: \textbf{¿qué te dijiste a ti mismo} para justificarlo? \emph{«No era para tanto», «después le digo», «no quería quedar mal».} \textbf{Tercera parte.} ¿Hubo alguien más ahí que \textbf{tampoco parecía muy convencido}? \emph{Piénsalo bien, porque casi siempre lo hay ---y casi nunca se sabe, porque nadie dice nada---.} \textbf{Y la última:} ¿alguna vez \textbf{tú} fuiste el de la presión ---el que insistió, el que dijo «no seas así»---?
\end{softbox}

{\sffamily\bfseries\fontsize{8}{10}\selectfont\textcolor{milcNegro!55}{MARCA CUANDO LO HAYAS HECHO}}
\milccheckrow{Escribiste tres síes que no querías decir, con cuánta gente había mirando.}
\milccheckrow{Anotaste la frase con que te lo justificaste a ti mismo.}
\milccheckrow{Pensaste si había alguien más que tampoco estaba convencido.}

\begin{infoband}{milcVerde}


\textbf{En tu cuaderno (Actividad 1 · IDENTIFICA).} Título: «Actividad 1. Los síes que no eran». \textbf{Formato:} los tres casos, con quién pidió, cuánta gente había y la frase con que te lo justificaste. \textbf{Extensión:} media página. \textbf{Sabes que terminaste cuando} identificaste \textbf{al menos una persona} que tampoco parecía convencida. Esta página es tuya: \textbf{nadie más tiene que leerla si tú no quieres}.
\end{infoband}

\puente{Ya los tienes. Ahora, por qué la unanimidad aprieta tanto y qué le hace un solo aliado.}

\milcestacion{milcTurquesa}{Fase 2 · Sistematización: entender la presión}

\begin{softbox}{milcTurquesa}{milcGris}{Cuatro claves para negarse}
Cuatro claves para que negarse deje de depender del valor que uno tenga ese día.
\end{softbox}

\milcpaso{1}{milcTurquesa}{\textbf{Casi nadie cree del todo lo que está diciendo.} En el experimento de Asch, la persona veía perfectamente cuál raya era ---la tarea era obvia--- y aun así repitió el error del grupo cerca de un tercio de las veces. \emph{No porque se hubiera convencido:} porque \textbf{quedar solo contra todos es muy caro}. \textbf{Ahora date vuelta y mira lo que eso significa para tu salón:} cuando el grupo entero parece estar de acuerdo en algo, \textbf{probablemente varios están haciendo exactamente lo que tú} ---seguir la corriente para no quedar solos---. \emph{La unanimidad casi siempre es falsa. Lo que es real es que nadie quiere ser el primero.}}
\milcpaso{2}{milcTurquesa}{\textbf{Un solo aliado desarma la presión.} Este es el hallazgo que hay que memorizar: cuando Asch puso \textbf{a una persona más} contestando distinto del grupo, la conformidad \textbf{se desplomó} ---y seguía funcionando aunque ese aliado también se equivocara, pero en otra dirección---. \textbf{La fuerza del grupo no está en el número: está en la unanimidad.} \emph{Rota la unanimidad, el poder se cae.} \textbf{La consecuencia práctica es enorme y casi nadie la usa:} no necesitas convencer a nadie ni ganar una discusión. \textbf{Necesitas una persona}, avisada de antemano, que diga «yo tampoco».}
\milcpaso{3}{milcTurquesa}{\textbf{El no se prepara antes, porque en el momento no hay tiempo.} Cuando llega la presión hay \textbf{dos o tres segundos} y diez personas mirando: nadie inventa una buena frase ahí. \emph{Por eso los «no» improvisados salen mal:} salen como pelea ---«¡ay, no jodás!»--- o como excusa que se puede desarmar ---«es que mi mamá no me deja»---, y a la excusa siempre le encuentran solución. \textbf{Un buen no tiene tres características:} es \textbf{corto}, \textbf{no pide permiso} y \textbf{no abre discusión}. \emph{«No, gracias» · «hoy no» · «paso» · «con eso no cuenten conmigo».} \textbf{Y algo que sirve más de lo que parece:} \emph{repetir la misma frase} en vez de buscar una nueva razón cada vez.}
\milcpaso{4}{milcTurquesa}{\textbf{Negarse a la cosa, no a la gente.} La mayoría no dice que no por miedo a la actividad, sino \textbf{por miedo a perder a la gente}. \emph{Por eso el mejor no separa las dos cosas:} rechaza \textbf{lo que se propone} y \textbf{deja intacto el vínculo}. \textbf{La forma más simple es ofrecer otra cosa:} «yo eso no, pero de una si vamos a\ldots». \emph{Y hay algo que conviene saber:} el grupo casi nunca expulsa a alguien por un no. \textbf{Lo que castiga es el sermón} ---negarse y encima explicar por qué los demás están mal---. \emph{Un no que no juzga a nadie casi siempre se acepta a la segunda.}}

\begin{drawbox}{milcTurquesa}{Las tres negativas que conviene tener listas}
\begin{pasoslista}
\item \textbf{La corta.} «No, gracias» · «hoy no» · «paso». \emph{Sin razón. Las razones se discuten; las frases no.}
\item \textbf{La que ofrece otra cosa.} «Eso no, pero de una si\ldots». \emph{Rechaza la cosa y conserva a la gente.}
\item \textbf{La que se apoya en un tercero.} «Nosotros dos pasamos». \emph{La más poderosa, porque rompe la unanimidad.}
\item \textbf{Y la regla de todas:} \textbf{repetir}, no argumentar. \emph{Quien discute su no ya lo puso en negociación.}
\end{pasoslista}
\end{drawbox}

\begin{softbox}{milcMostaza}{milcGris}{Errores comunes y su corrección}
\textbf{La excusa:} «es que mi mamá\ldots» ---a las excusas les encuentran solución---. \textbf{El sermón:} negarse y explicar por qué los demás están mal es lo único que sí se castiga. \textbf{Esperar a tener valor:} el valor no llega solo; llega con la frase preparada y con el aliado. \textbf{Irse sin decir nada:} evita el momento y deja el problema para la próxima. \textbf{Creer que estás solo:} en un grupo de veinte casi nunca hay veinte convencidos. \textbf{Discutirlo:} el que argumenta su no lo puso a negociar. \textbf{Buscar aliado en el momento:} se busca antes, en frío.
\end{softbox}

\begin{infoband}{milcTurquesa}


\textbf{En tu cuaderno (Actividad 2 · EXPLICA).} Título: «Actividad 2. La unanimidad». \textbf{Formato:} explica con tus palabras qué pasó en el experimento de Asch y qué cambió con el aliado; después, vuelve a uno de tus tres casos y escribe \textbf{quién habría podido ser tu aliado ahí}. \textbf{Extensión:} media página. \textbf{Sabes que terminaste cuando} puedes explicar por qué \textbf{la fuerza del grupo está en la unanimidad y no en el número}.
\end{infoband}

\puente{Ya sabes de dónde sale la fuerza. Es hora de armar tus tres negativas ---y de conseguir a tu aliado, que es la parte que de verdad funciona---.}

\milcestacion{milcMagenta}{Fase 3 · Praxis: armar el repertorio}

\begin{softbox}{milcMagenta}{milcGris}{Cómo se arma un no que no cuesta el puesto}
Ahora se arma el repertorio. Se escribe, se dice en voz alta ---que es donde se cae el que no ensayó--- y se consigue el aliado, que es la parte que de verdad funciona.
\end{softbox}

\milcpaso{1}{milcMagenta}{\textbf{Las tres negativas (10 min).} Escríbelas \textbf{con tus palabras}: la corta, la que ofrece otra cosa y la que se apoya en un tercero. \textbf{Condiciones:} que cada una quepa en \textbf{una frase}, que \textbf{no pida permiso} y que \textbf{no juzgue a nadie}. \emph{Léelas buscando la trampa más común: si alguna empieza con «es que», es una excusa y te la van a desarmar.}}
\milcpaso{2}{milcMagenta}{\textbf{El aliado (8 min).} Escoge a \textbf{una persona} de tu salón ---no tiene que ser tu mejor amigo; a veces funciona mejor que no lo sea--- y escribe \textbf{qué le vas a decir} para pedirle que sea tu aliado. \emph{Es más fácil de lo que parece, porque el trato es de ida y vuelta:} «si yo digo que no a algo, tú dices \emph{yo tampoco}; y yo hago lo mismo por ti». \textbf{Eso es todo. No hay que jurar nada ni armar un pacto solemne.}}
\milcpaso{3}{milcMagenta}{\textbf{El ensayo en voz alta (10 min, en tríos).} Uno presiona ---\textbf{con una situación inventada, nunca una real de alguien del salón}---, otro usa sus negativas y el tercero observa \textbf{cuánto tardó en decirla} y si sonó a no o a excusa. Roten. \emph{Después, la segunda ronda con la mejor parte:} \textbf{ahora el tercero hace de aliado} y dice «yo tampoco». \textbf{Van a sentir la diferencia en el cuerpo:} negarse solo es dificilísimo; de a dos, casi no cuesta.}
\milcpaso{4}{milcMagenta}{\textbf{La petición real (fuera de clase).} Háblale a la persona que escogiste y \textbf{hazle la propuesta de verdad}. \emph{Ese es el producto, y es lo único de esta guía que cambia algo mañana.} \textbf{Anota qué te dijo} ---y si dijo que no, anda donde otro: no es un rechazo, es que no todo el mundo quiere ese papel---.}

\begin{drawbox}{milcMagenta}{Antes de cerrar tu repertorio}
\begin{checklista}
\item Tus tres negativas caben \textbf{cada una en una frase}.
\item Ninguna empieza con «es que» ni pide permiso.
\item Ninguna juzga a los demás ni explica por qué están mal.
\item Las dijiste \textbf{en voz alta} y sonaron a no, no a excusa.
\item Escogiste a un aliado concreto y \textbf{ya le hiciste la propuesta}.
\item Entendiste que el trato es de ida y vuelta: también te toca decir «yo tampoco» por él.
\end{checklista}
\end{drawbox}

\begin{softbox}{milcMagenta}{milcGris}{Plantilla de trabajo}
\textbf{La corta.} \emph{«\_\_\_\_»} (una frase, sin razón, sin permiso). \\[1mm]
\textbf{La que ofrece otra cosa.} \emph{«Eso no, pero de una si \_\_\_\_.»} \\[1mm]
\textbf{La del aliado.} \emph{«Nosotros dos pasamos.»} ---y la propuesta: \emph{«si yo digo que no, dices \emph{yo tampoco}; y yo hago lo mismo por ti»}---.
\end{softbox}

\begin{infoband}{milcMagenta}


\textbf{En tu cuaderno (Actividad 3 · CREA).} Título: «Actividad 3. Mi repertorio del no». \textbf{Formato:} las tres negativas + el nombre del aliado y qué le dijiste + cómo salió el ensayo + qué respondió. \textbf{Extensión:} una página. \textbf{Sabes que terminaste cuando} \textbf{ya se lo pediste a alguien} ---no cuando lo tienes escrito---.
\end{infoband}

\puente{El producto es un repertorio ensayado. \emph{Un no que se improvisa casi siempre sale como pelea o como excusa; uno preparado sale como un no.}}

\milcestacion{milcMostaza}{Producto de la sesión}
\textbf{Repertorio del no: tres negativas ensayadas y un aliado al que ya se le pidió}.
\begin{softbox}{milcMostaza}{milcGris}{Criterios de calidad}
(1) Las tres negativas caben en una frase, no piden permiso y no abren discusión. (2) Ninguna funciona como excusa desarmable ni como sermón sobre los demás. (3) Se ensayaron en voz alta, con y sin aliado, y se registró la diferencia. (4) La petición al aliado está hecha de verdad, no solo escrita, y es recíproca. (5) Se explica por qué un solo aliado desarma la presión: rompe la unanimidad.
\end{softbox}

\puente{Antes de cerrar, revisa algo: si tus tres negativas suenan a discurso, no las vas a usar. Tienen que caber en una frase.}

\milcestacion{milcVino}{Evaluación liberadora · cinco dimensiones}

\begin{softbox}{milcVino}{milcGris}{Sentido de la evaluación}
Evaluar no es solo revisar una nota. Es reconocer qué aprendí, qué cambió en mí, qué emociones manejé, cómo me relaciono con la ciudad y la comunidad, y qué le dejaré a quien viene detrás.
\end{softbox}

\textbf{Mi reflexión liberadora en cinco dimensiones.} Completa cada frase iniciada:

\small
\begin{tabularx}{\textwidth}{>{\bfseries\RaggedRight\arraybackslash}p{3.4cm}Y>{\RaggedRight\arraybackslash}p{4.6cm}}
\rowcolor{milcVino}\color{white}Dimensión & \color{white}\bfseries Mi respuesta (frame starter) & \color{white}\bfseries Algo que descubrí\\
1. Personal & Lo que más cambió en mí fue \linead{6.5cm} & \linead{4.2cm}\\
2. Emocional & La emoción más fuerte fue \linead{2.5cm} y la regulé \linead{3cm} & \linead{4.2cm}\\
3. Ciudadana & En democracia esto importa porque \linead{6cm} & \linead{4.2cm}\\
4. Local (Cartago) & En mi barrio sería útil para \linead{6.5cm} & \linead{4.2cm}\\
5. Intergeneracional & A mi abuela/abuelo le diría \linead{6.5cm} & \linead{4.2cm}\\
\end{tabularx}
\normalsize

\begin{infoband}{milcVino}
\textbf{Para la próxima sustentación.} Vuelve a esta tabla en seis meses. Verás que las respuestas cambian — no porque lo aprendido cambie, sino porque tú cambias. La evaluación liberadora es una conversación contigo a través del tiempo.
\end{infoband}


\puente{Cerramos con tres voces: el que siempre es más que el grupo, lo que sí está en tu poder ---incluido el rechazo--- y qué pasa cuando el que dice que no queda registrado.}

\milcestacion{milcGradeDark}{Triángulo de pensamiento · cierre filosófico}

\begin{softbox}{milcVino}{milcGris}{Dussel · lente del nosotros}
\textit{«Analéctico quiere indicar el hecho real humano por el que todo hombre, todo grupo o pueblo se sitúa siempre más allá (aná-) del horizonte de la totalidad.»}

{\footnotesize\itshape\color{milcNegro!70}--- Enrique Dussel · Filosofía de la liberación (1977), §2.4}

Un grupo que presiona se comporta como si fuera \textbf{el mundo entero}: por eso quedar por fuera se siente como quedarse sin nada. \emph{Dussel dice que eso es falso:} cualquier persona ---y cualquier pueblo--- está \textbf{siempre más allá} del horizonte del grupo que lo rodea. \textbf{Los misak lo hicieron literal:} no discutieron dentro de las reglas del que los gobernaba, \textbf{se pararon en su propio Derecho Mayor} y desde ahí dijeron que no. \emph{Y ahí está la lección práctica:} un no es mucho más firme cuando \textbf{se apoya en algo tuyo} ---lo que quieres hacer, lo que ya te comprometiste a hacer, quién eres--- que cuando es solo un rechazo suelto.

\textbf{Cuando digo que no, ¿me estoy apoyando en algo mío, o solo estoy tratando de aguantar?}
\end{softbox}
\linead{\linewidth}\\[1mm]
\linead{\linewidth}

\begin{softbox}{milcMostaza}{milcGris}{Estoicismo · Epicteto · Enquiridión, 1 (c. 125 d.C.; trad. desde E. Carter) · cuidado interior}
\textit{«Algunas cosas están en nuestro poder y otras no. Están en nuestro poder la opinión, el impulso, el deseo, el rechazo; en una palabra, todo aquello que es obra nuestra.»}

{\footnotesize\itshape\color{milcNegro!70}--- Epicteto · Enquiridión, 1 (c. 125 d.C.; trad. desde E. Carter)}

Fíjate en la lista de Epicteto, porque una de las cuatro palabras es exactamente el tema de hoy: \textbf{el rechazo}. \emph{Está en tu poder ---no en el de ellos---.} \textbf{Lo que \emph{no} está en tu poder es cómo lo tomen}: si se enojan, si se ríen, si te dejan de hablar una semana. \emph{Y confundir las dos cosas es lo que paraliza:} uno no dice que no porque quiere controlar la reacción, \textbf{que es justamente lo único que no puede controlar}. \textbf{La versión útil:} \emph{tu tarea es decirlo bien y corto. La reacción es de ellos.}

\textbf{¿Estoy callándome porque no sé qué decir, o porque quiero controlar cómo van a reaccionar?}
\end{softbox}
\linead{\linewidth}\\[1mm]
\linead{\linewidth}

\begin{softbox}{milcTurquesa}{milcGris}{Floridi · lente de la infoesfera}
\textit{«Las tecnologías digitales están re-ontologizando la naturaleza misma de la infoesfera.»}

{\footnotesize\itshape\color{milcNegro!70}--- Luciano Floridi · La cuarta revolución (2014; trad.)}

Hoy la presión no se acaba cuando te vas para la casa: \textbf{sigue en el grupo, escrita}. \emph{Y eso tiene dos caras.} \textbf{La mala:} negarse queda registrado, y a veces con capturas ---la presión, que antes se disolvía, ahora se puede reenviar---. \textbf{La buena, y es la que hay que usar:} \emph{en un chat el aliado es más fácil que en el salón}. \textbf{No hay que pararse delante de nadie: basta con escribir «yo tampoco»}, y eso rompe la unanimidad igual de bien. \emph{Y hay una jugada todavía mejor, que casi nadie hace: escribirle al otro por privado antes ---«si dices que no, yo te apoyo»---. La fila deja de ser unánime antes de que empiece.}

\textbf{¿A quién del grupo podría escribirle por privado un «si dices que no, yo te apoyo»?}
\end{softbox}
\linead{\linewidth}\\[1mm]
\linead{\linewidth}

\begin{infoband}{milcNegro}
\textbf{Nota para el docente.} Las tres son citas textuales y llevan su referencia debajo. Puedes remitir a la obra en clase.
\end{infoband}

\milcestacion{milcVino}{Mi autoevaluación · marca una casilla por fila}

{\small
\setlength{\extrarowheight}{2.2pt}
\begin{tabularx}{\textwidth}{>{\RaggedRight\arraybackslash\bfseries}p{3.0cm}YYY}
\rowcolor{milcVino}\color{white}\bfseries Criterio & \color{white}\bfseries Lo logré & \color{white}\bfseries En proceso & \color{white}\bfseries Necesito apoyo\\[.6mm]
Comprensión del tema & \milcopcion{Explico las ideas con mis palabras.} & \milcopcion{Reconozco las ideas, falta precisión.} & \milcopcion{Necesito repasar lo básico.}\\[.8mm]
\rowcolor{milcGris}Calidad del repertorio & \milcopcion{Tengo tres negativas ensayadas y le pedí a alguien que me respalde.} & \milcopcion{Sé qué quiero decir, pero solo se me ocurre discutir o irme.} & \milcopcion{Sigo diciendo que sí y arrepintiéndome después.}\\[.8mm]
Aplicación y producto & \milcopcion{Mi producto cumple los criterios.} & \milcopcion{Está armado, con ajustes pendientes.} & \milcopcion{Aún está en construcción.}\\[.8mm]
\rowcolor{milcGris}Reflexión 5 dimensiones & \milcopcion{Respondí cada una con honestidad.} & \milcopcion{Respondí algunas con detalle.} & \milcopcion{Mis respuestas son muy generales.}\\[.8mm]
Triángulo de pensamiento & \milcopcion{Conecté las 3 voces con mi trabajo.} & \milcopcion{Conecté al menos una.} & \milcopcion{No las relaciono con mi proceso.}\\[.8mm]
\end{tabularx}}

\vspace{3mm}
\textbf{Hoy entendí que:}\\[1mm]
\linead{\linewidth}\\[1mm]
\linead{\linewidth}

\textbf{Mi compromiso para la próxima sesión es:}\\[1mm]
\linead{\linewidth}\\[1mm]
\linead{\linewidth}

\begin{softbox}{milcMostaza}{milcGris}{Cita de despedida · Marco Aurelio}
\textit{``El obstáculo en el camino se vuelve el camino. Lo que estorba al camino se vuelve camino.''}\\[1mm]
Cada error de hoy, bien observado, se convierte en pieza del camino que pisarás mañana.
\end{softbox}

\small
\begin{tabularx}{\textwidth}{>{\bfseries}p{3.6cm}Y}
\rowcolor{milcGradeDark!12}Nombre completo: & \linead{10cm}\\
Fecha: & \linead{4cm} \quad Curso/grupo: \linead{4cm}\\
Mi firma: & \linead{9cm}\\
\end{tabularx}
\normalsize

\begin{infoband}{milcGradeDark}
\textbf{Para la próxima sesión.} Dos cosas. \textbf{Primero, usa el repertorio una vez} ---aunque sea en algo mínimo--- y, sobre todo, \textbf{sé el aliado de alguien}: cuando veas a alguien negándose solo, di «yo tampoco». \emph{Anota qué pasó en los dos papeles, porque se sienten muy distinto.} \textbf{Segundo, pregunta en tu casa por un no difícil:} averigua con alguien mayor \textbf{una vez en que le tocó negarse a algo} ---en el trabajo, en el barrio, en la familia---: qué le pidieron, qué dijo, si alguien lo respaldó y qué pasó después. \textbf{Trae:} cómo te fue y lo que te contaron. \emph{Vas a encontrar dos cosas que se repiten: casi nadie recuerda el no como el momento heroico que uno imagina, y casi todos recuerdan si hubo alguien más de su lado.}
\end{infoband}

\end{document}
